\documentclass[12pt,a4paper]{article}
\usepackage[utf8]{inputenc}
\usepackage[french]{babel}
\usepackage{amsmath}
\usepackage{amsfonts}
\usepackage{amssymb}
\usepackage {hyperref} % couleur du tableau
\usepackage{xcolor}       % couleur du tableau
\usepackage{graphicx}
\usepackage{multirow} %Utilisation pour fusion case tableau%
\usepackage{colortbl} %COuleur Tableau%
\usepackage{fancyhdr}
\usepackage{times}
\usepackage[final]{pdfpages} 


\usepackage[left=4cm,right=3cm,top=2cm,bottom=2.5cm]{geometry} %FUCKING MARGE DE PETIT
\usepackage{setspace}
\setstretch{1,5} % FUCKING INTERLIGNE DE PETIT
\author{Lucas}
\sloppy

\title{Rapport de soutenance}

\renewcommand{\headrulewidth}{0.4pt}
\renewcommand{\footrulewidth}{0.4pt}
\fancyhead[L]{Info Spé}
\fancyhead[R]{Epita 2019}
\fancyfoot[L]{RTFPicture}

 \pagestyle{fancy}

\begin{document}

%Page de garde%


\begin{titlepage}


\flushright \includegraphics[scale=.15]{Pictures/EPITA_LOGO_DEF.png}
\flushright Promo 2019


%\includegraphics[scale=.15]{EPITA_LOGO_DEF.png}
%\hfill
%EPITA
%Promo 2019

\begin{center}

%\advance\leftskip by 1.75cm
\hspace{1.5cm}
{\Huge Rapport de 1er soutenance} 
\newline
\newline
\newline
\newline
\newline
\includegraphics[scale=1]{titre1.png}
%/ \includegraphics[scale=.35]{Pictures/logov2-0.jpg}
\end{center}
\begin{center}

%\vspace{0.7cm}
\end{center}
\begin{Large}
\begin{center}
Dan AZOULAY, Melvyn PETROCCHI, Ismail KESKAS, Lucas DESSERT \end{center}
\end{Large}
\vspace{.5cm}
\begin{center}
\begin{large}
Novembre 2015
\end{large}
\end{center}

\noindent
\newline
\newline
\textit{Nom du groupe :} Schrödinger \newline 
\textit{Nom du projet :} RTFPicture\newline 
\textit{Chef de projet :} Keskas Ismail \newline
\end{titlepage}
\newpage

%% ICI sommaire %%
\newpage
\tableofcontents

%% Debut Document%%
\newpage


\section*{Introduction}
Ce document est le rapport de la première soutenance. Il a pour but de présenter une synthèse sur le travail fourni par l'équipe en charge du projet. Cette équipe est formée de quatre étudiants en deuxième année du cycle préparatoire de l'EPITA : Dan AZOULAY (azoul\_d), Melvyn PETROCCHI (petroc\_m), Ismail KESKAS (keska\_i) et Lucas DESSERT (desser\_l).

 Notre projet est un logiciel de reconnaissance facial codé en C99 sous Linux, qui sera capable de rechercher des personnes présentes sur une photo donné, et qui à partir de cette photo donnera la liste des personnes présentes sur la photo qui auront été préalablement enregistré dans la base de donnée du logiciel.



\newpage


%/ $$$$$$$$$$$$$$$$$$$$$$$$$$$$$$$$$$$$$$$$$$$$$$$$$$$$$$$$$$$$$
 \section{Présentation du groupe Schrödinger} 
 


\subsection{Dan AZOULAY}


\subsection{Melvyn PETROCCHI}
Ce projet aura été une expérience professionnelle d'un atout
majeur en ce qui concerne le travail d'équipe. Je ne suis pas une
personne qui aime travailler en équipe au premier abord; je n'ai jamais réellement travaillé
en équipe autre que lorsque je sors en compétition d'équitation et lors de notre projet de SUP, un jeux video.
Mener une équipe avec un animal n'est pas du tout la même
ambiance que lorsque l'on mène une équipe avec des êtres
humains. Aujourd'hui, nous produisons un logiciel utilitaire ce qui est pour nous une nouvelle expérience.
Cette expérience est plutôt fructueuse, au point de vue où
cela nous apporte vraiment une façon de penser complètement
différente de celle que l'on aurait le réflexe d'avoir. 
Lorsque l'on mène une équipe
avec un animal, nous sommes à son écoute et sensible à tout ce
qu'il demande, tout en restant maître de lui en le dirigeant dans
une conance absolue. Cette complicité est aussi recherchée dans
le travail de groupe avec des humains également, mais celle-ci ne
se ressent pas de la même façon. Autre que les complicités
silencieuses que j'ai pu ressentir avec mes chevaux, en six mois,
j'ai appris à créer des liens de travail productifs et ecaces avec
d'autres personnes. Et cela, an de monter un projet qui soit au
nal comme nous l'avons prédit au début. Il me tarde de vivre des
expériences comme celle-ci an de savoir réellement à quoi cela
ressemble.
La première approche, le doute de l'ecacité me paraissait
assez absurde mais tout au long de nos travaux de partage et de
nos échanges je me suis aperçu que ça pouvait être vraiment
bénéque pour tout un projet que nous montons ensemble. Il
43
serait intéressant de se poser la question suivante : si une fois
arrivé à la n du projet serai-je capable de renouveler une telle
expérience ? Pour ma part, ce qui fut au nal pas si désagréable,
je pense que je pourrais facilement recommencer un projet en
équipe parce que cela nous permet d'avancer beaucoup plus vite.
Notamment, sur des problèmes posés par les résolutions diciles
nous avons eu souvent du mal à résoudre seul. Même si coder des
projets tout seul reste un plaisir, un projet en équipe de temps en
temps je pense serait la bienvenue.
En eet, cela nous permet d'apprendre diverses choses
permettant ainsi de progresser plus tard dans nos projets en solo.
Sachant que le métier d'ingénieur est un métier qui s'exerce
souvent en équipe, il me semble normal de le pratiquer
régulièrement, dans le but d'avoir des bases solides pour acquérir
des remèdes de conits ecaces.
Le métier de demain qui s'ouvre à nous ne sera entre autre
qu'un travail d'équipe.

\subsection{Ismail KESKAS}


\subsection{Lucas DESSERT}




\newpage
%/ $$$$$$$$$$$$$$$$$$$$$$$$$$$$$$$$$$$$$$$$$$$$$$$$$$$$$$$$$$$$$
\section{RTFPicture}
\subsection{Fonctionnement}

L'utilisation du logiciel est assez simple. Tous d'abord l'utilisateur devra remplir sa base de données de visages et nommé les personnes ajouté pour que le logiciel puisse ensuite traité ces photos, en extraire les informations qui lui seront utiles pour une futur reconnaissance, et la sauvegarder dans sa base de donnée.

Une fois ceci fait, l'utilisateur peux ensuite ajouter une photo quelconque au logiciel qui l'analysera et lui dira quels sont les personnes qu'il a reconnue en croisant les visages détectés sur la photo et sa base de données. Le logiciel sera aussi capable de signaler les absents, et proposera a l'utilisateur d'ajouter les nouveaux visages détecté (s'il y en à) dans sa base de donnée pour pouvoir les reconnaitre la prochaine fois.

L'utilisateur pourra a tous moment gérer la base de données, en ajoutant ou supprimant des personnes à sa guise (Il pourra bien sûr modifier le nom des personnes déjà enregistré s'il le souhaite).



\newpage
%/ $$$$$$$$$$$$$$$$$$$$$$$$$$$$$$$$$$$$$$$$$$$$$$$$$$$$$$$$$$$$$
\section{Détection de visage}



\subsection{Les caractéristique de Haar}



\subsection{Adaboost}


\newpage
%/ $$$$$$$$$$$$$$$$$$$$$$$$$$$$$$$$$$$$$$$$$$$$$$$$$$$$$$$$$$$$$
\section{Base de donnée}

\subsection{Sa structure}





 Affin de gérer principalement l'ajout et la suppression des visages dans le logiciel, il nous fallait créer une base de donné adapté a nos besoin. Elle doit pouvoir enregistrer une nouvelle personne, c'est a dire son nom, une ou plusieurs photos, et stocker les résultats du pré-traitement de l'image pour seront utilisés pour la reconnaissance de la personne.

Au départ nous avons réfléchis à comment créer cette base de donnée, et quoi utiliser pour y arriver. Nous avons étudier les possibilités de la bibliothèque SQL par exemple, qui est l'une des plus connus pour les bases de donnés, mais nous avons finalement choisi de n'utiliser aucune bibliothèque particulière et de coder nous même en C notre propre structure de donné. Nous n'étions pas convaincue par l'utilité de SQL pour notre projet et il nous a paru plus simple de créer notre propre type de donné affin de l'adapter parfaitement au besoin, et de la moduler plus facilement. De plus, aucuns d'entre nous n'avait des connaissances (même basique) dans la SQL.


Nous avons codé notre base de données sous forme de liste dynamique chaînée, pour ne pas avoir de problème de taille maximum ou d'allocation mémoire trop petite.Chaque personne ajouté se retrouve être un élément d'une liste chaîné avec un pointeur sur l’élément suivant de la liste (s'il existe). Chaque élément de cette liste dispose de trois string, l'un pour le nom de la personne qui est enregistré dans cet élément, et les deux autres seront interprété comme des pointeurs sur un dossier, le premier pointeur pointera sur le dossier qui contiendra le ou les images de la personne, et le deuxième pointera sur le dossier dans le quel les résultats des traitements de l'image seront enregistré (dans des fichier .text). L'avantage de cette structure est qu'au lancement de du logiciel, un simple parcourt d'un fichier texte qui contient tous les noms enregistré permettra de reconstituer la liste dynamique représentant tous les visage que le logiciel est capable de reconnaître.

\subsection{AU démarrage}

Mais pour réussis à reconstituer cette liste chaîné nous avons nommé les trois string de chaque élément d'une façon spécifique. Par exemple, si on enregistre une personne sous le nom de Xavier, la string du nom sera évidement "Xavier", la string qui fait office de pointeur sur le dossier qui contient la ou les photos de Xavier contiendra "Xavier-Image" et la troisième string qui permettra de trouver le dossier ou sont enregistré les résultats du pré-traitement de l'image contiendra "Xavier-Character". De cette façon un simple fichier text qui contient tous les noms des personnes enregistré (un registre) permettra via une fonction que nous avons codé, de reconstituer la liste dynamique tel qu'elle était lors de la dernière activité du logiciel. 

\subsection{L'ajout}

Pour ajouter un élément a notre liste chaîné, nous avons codé une fonction qui prend en paramètre la liste chaîné ainsi que le nom de la personne a ajouter et qui vas donc créer un nouvel élément, remplir les trois string en fonction du nom de la personne, et l'ajouter à la tête de liste (car l'ajout en fin de liste impliquerai de parcourir toute la liste avant de procéder à l'ajout). 


En plus de ça, la fonction d'ajout va écrire dans le fichier texte qui sert de registre, le nom de cette personne, affin qu'elle puisse être aussi ajouter dans la liste chaîné au prochain démarrage du logiciel.


Nous avons aussi codé une fonction qui permet d'ajouter plusieurs élément d'un coup, elle prend en argument la liste dynamique ainsi qu'un tableau de string.

\subsection{La suppression}

La fonction de suppression elle, recherche l’élément à supprimer au sein de la liste dynamique, le supprime si elle le trouve et efface son nom dans le registre ainsi que les deux dossiers qui contiennent les photos et les résultats du pré-traitement de l'image.

\subsection{La modification}

Notre fonction de modification permet pour le moment de modifier le nom d'une personne déjà enregistré dans la base de donné (elle modifie l'élément dans la liste dynamique ainsi que le nom écrit dans le registre).


\newpage
%/ $$$$$$$$$$$$$$$$$$$$$$$$$$$$$$$$$$$$$$$$$$$$$$$$$$$$$$$$$$$$$
\section{Les prochains objectifs}
\subsection{Reconnaissance de visage}

Pour la soutenance finale, nous prévoyons d'implémenter une interface graphique simple et intuitive qui permettra via de simple boutons, de visualiser les photos des personnes déjà enregistré, de supprimer ou d'ajouter des personnes dans la base de donnée, et bien évidemment de lancer une recherche de personne présente sur une photo.


\newpage
\subsection{L'interface graphique}


\newpage
\subsection{l'optimisation}


\newpage

\section*{Conclusion}



\end{document}

